\paragraph{Problem E.2.4 (Open in this incomplete return).}
The smooth identification of the shifted fivefold (E.16) remains.
Its exact group, integral homology and constancy under displacement
are established, and E.2.6 gives the actual marked splice
$(E\times T^4)\cup(T^4\times E)$, including its integral
Mayer--Vietoris extensions and middle torsion of order24.
E.2.9 proves that its natural stable framing bounds over the
aspherical group model. E.2.13--E.2.15 give an actual circle
fibration, its nine-dimensional fibre and separating-torus
return, and their complete integral homology and variation.

The earlier second-homotopy-module gap is now resolved.
E.2.16 determines the full $\mathbb Z\pi_1(E)$-module and its
complete derived peripheral map; E.2.17 computes the full
$\mathbb ZG$-module of the fivefold, including the central
intersection extension, and its restriction to the fibre.
E.2.18 additionally computes the actual normalized cusp's
Whitehead image and proves that the free module axes cannot
be realized by $S^2\times S^4$ or $S^2\times S^7$ connected
summands of $E$ or the circle fibre, respectively. These are
specific higher attaching constraints, not merely additive
invariants.

What remains is a complete smooth identification by an
applicable classification theorem with all required invariants
computed, or a direct comparison of complete framed handle
constructions. The known splice still contains the actual
regular-fibre complement $E$. Its remaining higher attaching
and smooth framing data have not yet been converted to such
a classification or finite handle identification. The current
geometric route preserves the boundary marking while deforming
the cusp correction to the unitless canonical-bundle plumbing;
its normal framings, triple-point compatibility and the entire
exterior gluing must be included before that route can settle
the smooth question. No smooth classification follows from
homology, the second-homotopy module, or a chain cancellation
alone.

